% ============================================================
%  Algebraic Limits of the Cl(6) Programme --- LIMITS / NO-GO PAPER
%  Todd M. Firestone
%  DRAFT --- not for circulation
%  Assembled from approved section drafts (§§1--7 + front matter).
%
%  FLAGS USED IN THIS FILE:
%    % NOTE:   --- placeholder or content to confirm at submission
%    % VERIFY: --- bibliographic detail to confirm (vol/page/year)
%    % ASM:    --- assembly decision left open
% ============================================================

\documentclass[12pt]{article}

% --- Encoding & fonts ---
\usepackage[utf8]{inputenc}
\usepackage[T1]{fontenc}
\IfFileExists{lmodern.sty}{\usepackage{lmodern}}{}
\usepackage{microtype}

% --- Page geometry (arXiv standard margins) ---
\usepackage[top=1in, bottom=1in, left=1.25in, right=1.25in]{geometry}

% --- Mathematics ---
\usepackage{amsmath}
\usepackage{amssymb}
\usepackage{amsthm}
\usepackage{mathtools}

% --- Tables ---
\usepackage{booktabs}
\usepackage{array}

% --- Citations: numbered, sorted ---
\usepackage[numbers,sort&compress]{natbib}
\bibliographystyle{unsrtnat}

% --- Color (needed for blue!60!black link-color mixing below) ---
\usepackage{xcolor}

% --- Hyperlinks (arXiv-safe) ---
\usepackage[
colorlinks=true,
linkcolor=blue!60!black,
citecolor=blue!60!black,
urlcolor=blue!60!black,
pdftitle={Algebraic Limits of the Cl(6) Programme},
pdfauthor={Todd M. Firestone}
]{hyperref}

% ============================================================
%  THEOREM ENVIRONMENTS  (section-prefixed; shared counter)
% ============================================================
\newtheorem{theorem}{Theorem}[section]
\newtheorem{lemma}[theorem]{Lemma}
\newtheorem{proposition}[theorem]{Proposition}
\newtheorem{corollary}[theorem]{Corollary}
\theoremstyle{remark}
\newtheorem{remark}[theorem]{Remark}

% ============================================================
%  MACROS
% ============================================================
\newcommand{\cl}{\mathrm{Cl}}
\newcommand{\spin}{\mathrm{Spin}}
\newcommand{\su}{\mathrm{SU}}
\newcommand{\UU}{\mathrm{U}}
\newcommand{\tr}{\mathrm{Tr}}
\newcommand{\R}{\mathbb{R}}
\newcommand{\Cplx}{\mathbb{C}}
\newcommand{\HH}{\mathbb{H}}
\newcommand{\OO}{\mathbb{O}}
\newcommand{\Dix}{\R\otimes\Cplx\otimes\HH\otimes\OO}
\newcommand{\grtensor}{\mathbin{\widehat{\otimes}}}  % graded tensor product
\newcommand{\half}{\tfrac{1}{2}}
\newcommand{\Comm}{\mathrm{Comm}}

% ============================================================
\begin{document}

\title{%
	\textbf{Algebraic Limits of the \texorpdfstring{$\cl(6)$}{Cl(6)}
		Programme:}\\[4pt]
	\textbf{A Schur Obstruction, the Forced Extension to
		\texorpdfstring{$\cl(10)$}{Cl(10)},}\\[4pt]
	\textbf{and Two No-Go Results for the Electroweak Sector}%
}
\author{Todd M.\ Firestone}
% NOTE: Target: Advances in Applied Clifford Algebras (short communication),
% with Phys. Lett. B brief report as alternative. Same venue as Stoica (2018),
% engaged directly in Sec.~\ref{sec:deriv}.
\date{}
\maketitle

% ============================================================
\begin{abstract}
% ============================================================
We delimit what the Clifford-algebra $\cl(6)$ approach to Standard-Model
structure can and cannot determine, establishing three results.  First, a
\emph{Schur obstruction}: any weak isospin $\su(2)$ commuting with color on the
eight-dimensional $\cl(6)$ ideal acts only on the color singlets, leaving the
quark triplet invariant.  Quarks therefore cannot be weak doublets within
$\cl(6)$; the result is exact and independent of the chosen embedding, and shows
that $\cl(6)$ is algebraically maximal for QCD.  Second, we exhibit the minimal
resolution as the graded tensor product
$\cl(10)=\cl(6)\grtensor\cl(4)$: weak isospin acts on the adjoined $\cl(4)$
factor, the Pati--Salam embedding
$\spin(10)\supset\su(4)\times\su(2)_L\times\su(2)_R$ restores the doublet
assignment, and all sixteen charges of one fermion generation are reproduced.
This construction is a structural re-expression of known Pati--Salam content;
its value is as an explicit bridge among the lattice $\cl(6)$, the
Furey--Gresnigt algebraic, and the Catterall K\"ahler--Dirac programmes.  Third,
we prove two no-go statements bounding the dynamics: the diagonal electroweak
grade decomposition is inert, equal to one-half at all couplings, and carries no
dynamical information; and no parameter of the Higgs potential follows from the
algebra, leaving the electroweak scale and Higgs mass undetermined.  We clarify,
rather than refute, Stoica's ``broken by geometry'' characterization: the
geometry fixes the breaking \emph{pattern} and embedding ratios, not the
breaking \emph{dynamics}.
\end{abstract}

\bigskip
\noindent\textit{Keywords:} Clifford algebra; $\cl(6)$; Schur's lemma;
Pati--Salam; $\cl(10)$; lattice gauge theory; electroweak symmetry breaking;
derivability limit; no-go theorem.

% ============================================================
\section{Introduction}
\label{sec:intro}
% ============================================================

The Clifford algebra $\cl(6)$ has proven a natural algebraic home for the color
sector of the Standard Model.  Its eight-dimensional minimal left ideal carries
the $\su(3)_C$ representation content
$\mathbf{3}\oplus\bar{\mathbf{3}}\oplus\mathbf{1}\oplus\mathbf{1}$, and Papers~1
and~2 of this series showed that this structure supports genuine dynamical
calculation: Paper~1~\cite{FirestoneP1} verified a Wilson--Dirac operator built
in $\cl(6)$ against sixteen independent benchmarks, and
Paper~2~\cite{FirestoneP2} constructed and measured grade-decomposed QCD
correlators on quenched lattice ensembles.  The natural next question is whether
the rest of the Standard-Model gauge structure---in particular weak isospin
$\su(2)_L$---can be accommodated within the same algebraic setting.

This paper answers that question by establishing where the algebraic machinery
reaches and where it stops.  We organize the paper around three results.  First
(Sec.~\ref{sec:schur}), a Schur obstruction: no $\su(2)$ commuting with color
can act on the $\cl(6)$ ideal except on the color singlets, so quarks cannot be
weak doublets within $\cl(6)$.  The result is exact and embedding-independent,
and it shows that $\cl(6)$ is algebraically maximal for QCD.  Second
(Sec.~\ref{sec:cl10}), the minimal resolution: the graded tensor product
$\cl(10)=\cl(6)\grtensor\cl(4)$ adjoins a factor on which weak isospin can act,
restoring the doublet assignment through the Pati--Salam embedding and
reproducing one full fermion generation.  Third (Secs.~\ref{sec:rigidity}
and~\ref{sec:deriv}), two no-go statements that bound the dynamics: the diagonal
electroweak grade decomposition is inert at all couplings
(Sec.~\ref{sec:rigidity}), and no parameter of the Higgs potential follows from
the algebra (Sec.~\ref{sec:deriv}).

These results sit alongside, and sharpen, an existing body of work.  Furey and
collaborators~\cite{Furey2012,Furey2014,Furey2015,Furey2018a,Furey2018b} have
developed the algebraic Standard-Model programme from division algebras,
recovering the gauge group and one generation's quantum numbers;
Stoica~\cite{Stoica2018}, in this journal, constructed the electroweak sector in
the $\cl(6)$ language and characterized the electroweak symmetry as ``already
broken by the geometry of the algebra.''  We engage that characterization
directly in Sec.~\ref{sec:deriv}, where we argue it should be read as a
statement about the breaking \emph{pattern} and embedding ratios rather than the
breaking \emph{dynamics}---a clarification, not a contradiction.

The paper is therefore a limits note within the series: its contribution is to
delimit precisely what the $\cl(6)$ construction determines, to exhibit the
forced extension that lifts its central obstruction, and to bridge the lattice
construction of Papers~1--2 to the broader algebraic and lattice programmes it
neighbors.

% ============================================================
\section{The Schur Obstruction}
\label{sec:schur}
% ============================================================

\subsection{Setup}
\label{sec:schur-setup}

Papers~1 and~2 of this series established that the eight-dimensional minimal
left ideal of the Euclidean Clifford algebra $\cl(6)$, with generators
satisfying $\{e_i,e_j\}=2\delta_{ij}$, is the algebraic home of the QCD color
sector.  The color gauge group $\su(3)_C$ is embedded through the chain
$\spin(6)\supset\su(3)\times\UU(1)$, and under $\su(3)_C$ the ideal decomposes
into irreducible representations as
\begin{equation}
	\mathbf{8} \;=\; \mathbf{3}\ \oplus\ \bar{\mathbf{3}}\ \oplus\
	\mathbf{1}_a\ \oplus\ \mathbf{1}_b.
	\label{eq:decomp}
\end{equation}
The triplet $\mathbf{3}$ carries the quark color charge; the antitriplet
$\bar{\mathbf{3}}$ carries the antiquark charge; and the two color singlets
$\mathbf{1}_a,\mathbf{1}_b$ complete the space.  This decomposition is fixed by
the embedding and is verified at machine precision in Paper~1.

The question this section settles is whether a second non-abelian gauge factor
---specifically a weak-isospin $\su(2)_L$ commuting with color, as the Standard
Model requires---can be realized on this same eight-dimensional space.  The
answer is no, and the obstruction is exact and embedding-independent.

\subsection{The Commutant of \texorpdfstring{$\su(3)_C$}{SU(3)\_C}}
\label{sec:commutant}

In the Standard Model, color and weak isospin commute:
$[\su(3)_C,\su(2)_L]=0$.  Any candidate $\su(2)_L$ acting on the ideal must
therefore have all of its generators in the commutant of the $\su(3)_C$
representation.  By Schur's lemma~\cite{FultonHarris}, the commutant of a
completely reducible representation is determined entirely by its irreducible
content and the multiplicities with which the irreducibles appear.

Reading off the multiplicities from~\eqref{eq:decomp}:
\begin{itemize}
	\item The triplet $\mathbf{3}$ appears once.
	\item The antitriplet $\bar{\mathbf{3}}$ appears once.  For $\su(3)$ the
	fundamental and the antifundamental are \emph{inequivalent} irreducible
	representations---the fundamental is complex, not self-conjugate---so no
	nonzero intertwiner maps the $\mathbf{3}$ block to the $\bar{\mathbf{3}}$
	block.
	\item The trivial representation appears twice, as $\mathbf{1}_a$ and
	$\mathbf{1}_b$.  These two copies are equivalent (both trivial), so
	intertwiners may freely mix them.
\end{itemize}
Schur's lemma then gives the commutant as a direct sum of one matrix algebra per
equivalence class of irreducible, of size equal to the multiplicity:
\begin{equation}
	\Comm\bigl(\su(3)_C\bigr) \;=\; \Cplx\ \oplus\ \Cplx\ \oplus\ M_2(\Cplx),
	\label{eq:commutant}
\end{equation}
of total complex dimension $1+1+4=6$.  The first two summands are the scalar
multiples of the identity on the $\mathbf{3}$ and $\bar{\mathbf{3}}$ blocks
respectively; the third is the full algebra of $2\times2$ complex matrices
acting on the two-dimensional space spanned by the color singlets
$\{\mathbf{1}_a,\mathbf{1}_b\}$.

\subsection{The Obstruction}
\label{sec:obstruction}

The Lie algebra of the commutant~\eqref{eq:commutant} is
\begin{equation}
	\mathfrak{u}(1)\ \oplus\ \mathfrak{u}(1)\ \oplus\ \mathfrak{u}(2)
	\;=\; \mathfrak{u}(1)^{\oplus 3}\ \oplus\ \mathfrak{su}(2).
	\label{eq:liealg}
\end{equation}
Its only non-abelian simple ideal is the $\mathfrak{su}(2)$ contained in
$\mathfrak{u}(2)$.  That $\mathfrak{su}(2)$ acts on the two-dimensional space of
color singlets and is trivial on every other block.  We record the consequence
as a theorem.

\begin{theorem}[Schur Obstruction]
	\label{thm:schur}
	Let $\su(3)_C$ act on the eight-dimensional $\cl(6)$ ideal by the
	decomposition~\eqref{eq:decomp}.  Then every $\su(2)$ subgroup commuting
	with $\su(3)_C$ acts nontrivially only on the two color singlets
	$\mathbf{1}_a,\mathbf{1}_b$, and acts as the identity on the color triplet
	$\mathbf{3}$ and antitriplet $\bar{\mathbf{3}}$.  Consequently the quark
	triplet is a singlet under any such $\su(2)$, the exact opposite of the
	Standard Model assignment in which left-handed quarks form weak-isospin
	doublets.  The conclusion depends only on the irreducible content and
	multiplicities in~\eqref{eq:decomp}, and is therefore independent of the
	choice of generator matrices.
\end{theorem}

\begin{proof}
	By the commutativity requirement, the generators of any color-commuting
	$\su(2)$ lie in $\Comm(\su(3)_C)$.  By Schur's lemma this commutant
	is~\eqref{eq:commutant}, whose Lie algebra~\eqref{eq:liealg} has
	$\mathfrak{su}(2)$ as its only non-abelian simple ideal.  That
	$\mathfrak{su}(2)$ is supported on the span of
	$\{\mathbf{1}_a,\mathbf{1}_b\}$ and annihilates the $\mathbf{3}$ and
	$\bar{\mathbf{3}}$ blocks.  Hence any color-commuting $\su(2)$ leaves the
	triplet invariant.  The argument uses only~\eqref{eq:decomp}, so it holds
	for every embedding producing that decomposition.
\end{proof}

The inequivalence of $\mathbf{3}$ and $\bar{\mathbf{3}}$ is the load-bearing
step: it is what reduces each of those two blocks to a single scalar rather than
a larger matrix algebra.  Were the representation real or pseudoreal in those
blocks, the commutant would be larger and the conclusion could differ; for
$\su(3)$ it is not.

\subsection{Algebraic Maximality}
\label{sec:maximality}

\begin{corollary}[Algebraic Maximality for QCD]
	\label{cor:maximal}
	The $\cl(6)$ Wilson--Dirac framework of Papers~1 and~2 is algebraically
	maximal for QCD: $\su(3)_C$ exhausts the nontrivial gauge structure the
	eight-dimensional ideal can carry while keeping quarks in the color
	triplet.  A nontrivial weak-isospin $\su(2)_L$ acting correctly on quarks
	cannot be added within $\cl(6)$; it requires enlarging the algebra.
\end{corollary}

This reframes what might look like a limitation as a positive structural
statement.  $\cl(6)$ is exactly the right size for the color sector---neither
too small to carry $\su(3)_C$ nor large enough to accommodate a spurious extra
non-abelian factor on the quarks.  The next section shows that the minimal
enlargement consistent with the Standard Model assignment is forced to be
$\cl(10)$.

\subsection{Numerical Confirmation}
\label{sec:schur-numeric}

The theorem is exact algebra and needs no numerics, but a direct check rules out
the most natural attempt to build a color-commuting weak generator by hand.
Constructing a candidate $T_3$ in the eight-dimensional space and testing it
against the off-diagonal Gell-Mann generators of $\su(3)_C$ gives
\begin{equation}
	\bigl\|\,[\,T_3,\,T_a\,]\,\bigr\| \;=\; 1.0
	\qquad \text{for } a\in\{1,2,6,7\},
	\label{eq:commnumeric}
\end{equation}
i.e.\ the commutator fails to vanish.  This matches the prediction of
Theorem~\ref{thm:schur} that no such color-commuting $\su(2)$ generator exists
on the triplet.  The full check is reproduced in the verification script
accompanying this paper.

% ============================================================
\section{The Forced Extension:
	\texorpdfstring{$\cl(10)=\cl(6)\grtensor\cl(4)$}{Cl(10)=Cl(6)(x)Cl(4)}}
\label{sec:cl10}
% ============================================================

Theorem~\ref{thm:schur} showed that weak isospin has no color-commuting home
inside $\cl(6)$: any $\su(2)$ commuting with $\su(3)_C$ is supported on the color
singlets and leaves quarks invariant.  The Standard Model requires the
opposite---left-handed quarks are weak doublets.  The resolution is not a
different embedding of $\su(2)_L$ within $\cl(6)$, which Theorem~\ref{thm:schur}
rules out for all embeddings, but an enlargement of the algebra itself.  This
section exhibits the minimal such enlargement that restores the doublet
assignment and shows precisely how it sidesteps the obstruction.

\subsection{The Graded Tensor Product Construction}
\label{sec:graded-tensor}

We construct $\cl(10)$ as the graded ($\mathbb{Z}_2$-graded, or ``super'')
tensor product of $\cl(6)$ with $\cl(4)$,
\begin{equation}
	\cl(10) \;=\; \cl(6)\grtensor\cl(4),
	\label{eq:cl10}
\end{equation}
acting on the $32$-dimensional space $\Cplx^8\otimes\Cplx^4$.  Let
$\Gamma_7=i\,e_1e_2\cdots e_6$ denote the $\cl(6)$ chirality element, normalized
so that $\Gamma_7^2=+I_8$, and let $\{f_1,f_2,f_3,f_4\}$ generate $\cl(4)$ with
$\{f_j,f_k\}=2\delta_{jk}$.  The ten generators of $\cl(10)$ are
\begin{align}
	\Gamma_i     &= e_i\otimes I_4        && (i=1,\ldots,6),
	\label{eq:gen-a}\\
	\Gamma_{6+j} &= \Gamma_7\otimes f_j   && (j=1,\ldots,4).
	\label{eq:gen-b}
\end{align}
The factor $\Gamma_7$ in~\eqref{eq:gen-b} is the grading cocycle of the graded
tensor product.  It is exactly what is needed to make the cross anticommutators
vanish: for $i\in\{1,\ldots,6\}$ and $j\in\{1,\ldots,4\}$,
\begin{equation}
	\{\Gamma_i,\Gamma_{6+j}\}
	= \bigl(e_i\Gamma_7 + \Gamma_7 e_i\bigr)\otimes f_j
	= 0,
	\label{eq:cross}
\end{equation}
because $e_i$ anticommutes with the top-grade element $\Gamma_7$ in $\cl(6)$.
Without the $\Gamma_7$ insertion these would not vanish and~\eqref{eq:cl10}
would fail.  With it, one verifies the full Clifford relation
\begin{equation}
	\{\Gamma_a,\Gamma_b\} = 2\delta_{ab}\,I_{32},
	\qquad a,b\in\{1,\ldots,10\},
	\label{eq:closure}
\end{equation}
confirmed at machine precision in the verification script accompanying this
paper.

The factorization $\cl(6)\otimes\cl(4)\cong\cl(10)$, with the color structure
carried by the $\cl(6)$ factor and the weak structure by the $\cl(4)$ factor, is
due to Gresnigt~\cite{Gresnigt2020}; we adopt it directly.  Our contribution in
this section is the explicit \emph{graded} tensor-product realization with the
$\Gamma_7$ cocycle~\eqref{eq:gen-b}, which makes the $\cl(10)$ closure manifest,
and the use of that realization to lift the obstruction of
Theorem~\ref{thm:schur}; the underlying $\cl(6)\otimes\cl(4)$ assignment is
Gresnigt's.

\subsection{Chirality and the 16-Dimensional Weyl Spinor}
\label{sec:chirality}

The $\cl(10)$ chirality element factorizes cleanly across the two tensor
factors:
\begin{equation}
	\Gamma_{\mathrm{chiral}}^{\cl(10)}
	\;=\; \Gamma_7^{\cl(6)}\otimes\gamma_5^{\cl(4)},
	\label{eq:chiral-fact}
\end{equation}
where $\gamma_5^{\cl(4)}=f_1f_2f_3f_4$ is the $\cl(4)$ chirality.  Its $+1$
eigenspace is $16$-dimensional and carries exactly one Standard Model
generation.

\subsection{The Pati--Salam Embedding}
\label{sec:patisalam}

The $\spin(10)$ Lie algebra is spanned by the bivectors
$\Sigma_{ab}=\tfrac{i}{4}[\Gamma_a,\Gamma_b]$.  It contains the Pati--Salam
subalgebra
\begin{equation}
	\spin(10)\;\supset\;\su(4)\times\su(2)_L\times\su(2)_R,
	\label{eq:ps-embed}
\end{equation}
under which the $16$-dimensional Weyl spinor branches as
\begin{equation}
	\mathbf{16} \;=\; (\mathbf{4},\mathbf{2},\mathbf{1})\ \oplus\
	(\bar{\mathbf{4}},\mathbf{1},\mathbf{2}).
	\label{eq:branching}
\end{equation}
The $(\mathbf{4},\mathbf{2},\mathbf{1})$ carries the left-handed fields; the
$(\bar{\mathbf{4}},\mathbf{1},\mathbf{2})$ carries the right-handed fields as a
charge-conjugate multiplet.  This embedding and branching are
standard~\cite{PatiSalam1974,BaezHuerta2010}; we use them, and do not claim them
as new.  What is constructed here is their explicit realization through the
graded tensor product~\eqref{eq:cl10}--\eqref{eq:gen-b}.

\subsection{How the Enlargement Evades the Schur Obstruction}
\label{sec:evasion}

This is the structural point of the section, and it connects directly to
Theorem~\ref{thm:schur}.  Inside $\su(4)$ the color group sits as
\begin{equation}
	\su(4)\;\supset\;\su(3)_C\times\UU(1)_{B-L},
	\label{eq:su4}
\end{equation}
with $\su(3)_C$ residing entirely in the $\cl(6)$ tensor factor: the fundamental
$\mathbf{4}$ of $\su(4)$ is the $\Gamma_7=+1$ sector of the $\cl(6)$
construction, and the $B-L$ charge is proportional to the $\cl(6)$ Cartan
generator, $B-L=\tfrac{2}{3}H$.

Weak isospin $\su(2)_L$, by contrast, acts on the adjoined $\cl(4)$ factor---a
different tensor slot from color.  This is precisely why the obstruction of
Theorem~\ref{thm:schur} does not apply.  That theorem constrained any $\su(2)$
lying in the commutant of $\su(3)_C$ \emph{within} $\cl(6)$, forcing it onto the
color singlets.  Here $\su(2)_L$ does not live in $\cl(6)$ at all; it lives in
the $\cl(4)$ factor that the tensor product adjoins.  Because color acts as
$I_4$ on that factor~\eqref{eq:gen-a}, $\su(2)_L$ commutes with $\su(3)_C$
automatically and acts on the full $\su(4)$ multiplet---and therefore on the
quarks---as a genuine doublet symmetry.  The tensor enlargement is exactly the
minimal step that supplies weak isospin a home where it can act on quarks the
way the Standard Model demands.

\subsection{Charge Assignments}
\label{sec:charges}

The electroweak hypercharge and electric charge are recovered from the
Pati--Salam data by the standard relations
\begin{equation}
	Y = T_{3R} + \half(B-L),
	\qquad
	Q = T_{3L} + Y,
	\label{eq:YQ}
\end{equation}
with $T_{3L},T_{3R}$ the Cartan generators of $\su(2)_L,\su(2)_R$.  Evaluating
\eqref{eq:YQ} on all sixteen states of one generation---the left-handed doublets
$(\nu_L,e_L)$, $(u_L,d_L)$ and the right-handed singlets $e_R,u_R,d_R,\nu_R$---
reproduces every Standard Model charge, with residual at most
$1.4\times10^{-16}$ and total charge $\sum Q=0$ over the generation (verification
script, $16$-charge table).

One subtlety must be stated.  The right-handed fields sit in the
charge-conjugate block $(\bar{\mathbf{4}},\mathbf{1},\mathbf{2})$.  Consequently
$\su(2)_R$ must be represented on that block by $-\tau^*/2$ rather than
$+\tau/2$; representing it without the conjugation yields incorrect right-handed
charges.  This is a representation-theoretic bookkeeping point about the
conjugate multiplet, not a physical input, and it is built into the
construction.

\subsection{Status of the Result}
\label{sec:cl10-status}

We state plainly what this section is and is not.  The Pati--Salam content of
$\spin(10)$ is known and credited above; we make no claim of new physics here,
and by the informational-novelty standard applied throughout this series this
section is a structural re-expression rather than an informational advance.  Its
value is constructive and connective: it provides an explicit graded-tensor-
product realization $\cl(10)=\cl(6)\grtensor\cl(4)$ that (i)~takes the
lattice-ready $\cl(6)$ framework of Papers~1--2 as a literal tensor factor,
(ii)~restores the weak-doublet assignment that Theorem~\ref{thm:schur} forbids
inside $\cl(6)$ alone, and (iii)~thereby links three otherwise-separate
programmes---the lattice $\cl(6)$ construction of this series, the algebraic
Standard-Model work of Furey~\cite{Furey2018a,Furey2018b} and
Gresnigt~\cite{Gresnigt2020}, and the K\"ahler--Dirac\,/\,Pati--Salam lattice
work of Catterall~\cite{Catterall2023,Catterall2024}.  The extension is ``minimal'' in the
precise sense that it is the smallest graded-tensor enlargement of $\cl(6)$
restoring the doublet assignment; we do not claim it is the unique conceivable
enlargement.

% ============================================================
\section{Algebraic Rigidity of the
	\texorpdfstring{$T_3$}{T3}-Diagonal Grade Decomposition}
\label{sec:rigidity}
% ============================================================

Section~\ref{sec:deriv} will show that the algebra does not constrain the Higgs
potential.  This section establishes a companion negative result on the gauge
side: the diagonal part of the electroweak grade decomposition is dynamically
inert.  It equals one-half identically---for every gauge configuration and at
every coupling---and therefore carries no dynamical information.  We state this
precisely and prove it, then contrast it with the genuinely dynamical color
decomposition of Paper~2, since the two look superficially similar and are not
the same kind of object.

\subsection{The Decomposition}
\label{sec:rigidity-decomp}

Working in $\su(2)$ pure gauge theory implemented in the eight-dimensional space
of Paper~3's successor construction, the weak-isospin generators are embedded in
the second tensor factor, $T_a=I_2\otimes(\sigma_a/2)\otimes I_2$, and a gauge
link $U\in\su(2)$ is embedded as $I_2\otimes U\otimes I_2$.  The $T_3$ spectral
projectors are
\begin{equation}
	P_{\pm}^{\mathrm{EW}} \;=\; I_2\otimes P_{\uparrow/\downarrow}\otimes I_2,
	\label{eq:Ppm}
\end{equation}
with $P_{\uparrow}=(I+\sigma_z)/2$, $P_{\downarrow}=(I-\sigma_z)/2$.  Inserting
these into the gauge correlator splits it into diagonal and off-diagonal parts,
\begin{equation}
	G(t) \;=\; G^{(+)}(t) + G^{(-)}(t) + \Delta_{\mathrm{cross}}(t),
	\label{eq:Gdecomp}
\end{equation}
where $G^{(\pm)}(t)$ are the $T_3$-diagonal projected correlators and
$\Delta_{\mathrm{cross}}(t)$ is the off-diagonal contribution.  The off-diagonal
term is the carrier of the non-trivial dynamics; it reduces to
$\Delta_{\mathrm{cross}}(t)=8\sum_{x,\mu}\mathrm{Re}\,[\,b(x,t)\,\bar
b(x,0)\,]$, where $b$ is the off-diagonal $\su(2)$ link element.  We do not
develop $\Delta_{\mathrm{cross}}$ in this paper; it is the subject of later work
and is noted again in Sec.~\ref{sec:discussion}.

\subsection{The Rigidity Result}
\label{sec:rigidity-prop}

Define the diagonal grade fractions
\begin{equation}
	f_{\pm}(t) \;=\; \frac{G^{(\pm)}(t)}{G^{(+)}(t)+G^{(-)}(t)}.
	\label{eq:fpm}
\end{equation}

\begin{proposition}[Algebraic Rigidity of the Diagonal Decomposition]
	\label{prop:rigidity}
	For every gauge configuration, and at all values of the coupling,
	$f_{+}(t)=f_{-}(t)=\half$.  The diagonal grade fractions are constant by
	construction and contain no dynamical information.
\end{proposition}

\begin{proof}
	The diagonal projectors~\eqref{eq:Ppm} extract the $(1,1)$ and $(2,2)$
	matrix elements of the $\su(2)$-valued line product $W(t)=\prod U$ entering
	the correlator.  Since $W(t)\in\su(2)$, unitarity and unimodularity give
	$W_{22}=(W_{11})^{*}$, hence $|W_{11}|=|W_{22}|$.  The diagonal correlators
	$G^{(\pm)}$ depend on these matrix elements only through their moduli and do
	so symmetrically, so $G^{(+)}(t)=G^{(-)}(t)$ configuration by
	configuration, independent of the coupling.  The fractions~\eqref{eq:fpm}
	are therefore each $\half$.  The full tensor-index proof, including the
	rank-one collapse step, is given in Paper~4 of this series.
\end{proof}

The content of Proposition~\ref{prop:rigidity} is that the equality is exact and
unconditional.  There are no deviations from $\half$ to measure: the diagonal
decomposition is rigid.

\subsection{What This Result Is, and What It Is Not}
\label{sec:rigidity-contrast}

Proposition~\ref{prop:rigidity} resembles the color-sector result of Paper~2
only on the surface.  The two differ in kind, and conflating them would
misrepresent both.

In Paper~2 the color grade fractions satisfy $f_r=\dim(r)/8$ \emph{exactly in
the free field} and only \emph{approximately} in the interacting theory; the
deviations are genuine and dynamical.  The mechanism is traceable: the Cartan
grading operator is $H=B_{12}+B_{34}+B_{56}$, and the spacetime Dirac generators
$e_1,\ldots,e_4$ act on the same tensor factors (qubits~1 and~2) as $B_{12}$ and
$B_{34}$.  Dirac\,/\,chirality structure therefore feeds the color-grade label,
and the interacting-theory deviations record it.  That is a kinematic limit
theorem with measurable config-by-config departures.

In the present case the diagonal fractions equal $\half$ at \emph{all}
couplings, exactly.  This is an algebraic identity, not a kinematic limit
theorem, and it must not be called one.  The diagonal projection is sensitive
only to $|W_{11}|=|W_{22}|$; it is blind to the gauge dynamics by construction.
Where the color decomposition has a dynamical signal sitting on top of its
free-field value, the diagonal electroweak decomposition has nothing on top of
its constant value at all.

We also fix terminology.  The object in~\eqref{eq:Gdecomp} is the $T_3$-diagonal
grade decomposition of the $\su(2)$ gauge-link correlator.  It is not a
``$W$-boson correlator'': there is no $W$ boson in pure $\su(2)$ gauge theory
absent electroweak symmetry breaking, and no Higgs sector is present here.  And,
by the above, it is not governed by a kinematic limit theorem.  These naming
points are not pedantry; the earlier ``kinematic theorem'' framing obscured
exactly the free-versus-all-couplings distinction that distinguishes
Proposition~\ref{prop:rigidity} from the Paper~2 result.

\subsection{The No-Go Content}
\label{sec:rigidity-nogo}

Proposition~\ref{prop:rigidity} is the electroweak companion to the Derivability
Limit of Sec.~\ref{sec:deriv}.  Both delimit the algebraic machinery from below:
the algebra fixes neither the Higgs potential (Sec.~\ref{sec:deriv}) nor any
dynamics in the diagonal electroweak grade channel (this section).  The diagonal
grade decomposition is a consistency-bookkeeping device, not a probe of gauge
dynamics.  The channel that does carry dynamics is the off-diagonal
$\Delta_{\mathrm{cross}}$ of~\eqref{eq:Gdecomp}, which lies outside the scope of
this limits paper and is taken up separately (Sec.~\ref{sec:discussion}).

% ============================================================
\section{The Derivability Limit}
\label{sec:deriv}
% ============================================================

\subsection{Statement and Proof}
\label{sec:deriv-thm}

We state the central negative result on the scalar sector.

\begin{theorem}[Derivability Limit]
	\label{thm:deriv}
	No condition on the Higgs potential
	\begin{equation}
		V(\phi) = \mu^2\,|\phi|^2 + \lambda\,|\phi|^4
		\label{eq:V}
	\end{equation}
	follows from the algebraic structure of $\Dix$ alone.
\end{theorem}

\begin{proof}
	The algebraic structure encodes three things: the representation content of
	the fundamental fields, the action of the gauge group on those
	representations, and the charge assignments and symmetry relations among
	them.  It does not encode a choice of Lagrangian.

	The Higgs field transforms as the doublet $(\mathbf{2},\half)$ of
	$\su(2)_L\times\UU(1)_Y$.  For a single such doublet, $\phi^\dagger\phi$ is
	the unique quadratic gauge invariant, and every positive power of it is
	likewise invariant.  Hence for any $\mu^2\in\R$ and any $\lambda>0$, the
	potential~\eqref{eq:V} is fully invariant under $\su(2)_L\times\UU(1)_Y$:
	the algebraic structure is consistent with this choice of dynamics.

	Since the algebra is consistent with the entire two-parameter family
	$\{(\mu^2,\lambda):\mu^2\in\R,\,\lambda>0\}$, it constrains no member of it.
	In particular it is consistent with $\mu^2>0$ (unbroken symmetry),
	$\mu^2=0$ (massless tree-level scalar), and $\mu^2<0$ (the physical broken
	case), and with any magnitude of $\lambda$.  The algebra constrains neither
	the sign of $\mu^2$ nor the value of $\lambda$.
\end{proof}

\begin{corollary}
	\label{cor:vscale}
	The electroweak scale $v=\sqrt{-\mu^2/\lambda}$ is unconstrained by the
	algebra; in particular the algebra does not explain why $v$ is of order
	$246~\mathrm{GeV}$ rather than the Planck scale.
\end{corollary}

\begin{corollary}
	\label{cor:mhmass}
	The Higgs boson mass $m_H^2=2\lambda v^2=-2\mu^2$ is unconstrained by the
	algebra.
\end{corollary}

\subsection{What the Theorem Is: A Precise Statement of Standard Lore}
\label{sec:deriv-debill}

We are deliberate about the status of Theorem~\ref{thm:deriv}, because it is
easy to overstate.  At its core it is the standard effective-field-theory fact
that a symmetry determines the \emph{form} of the operators it permits, not their
\emph{coefficients}~\cite{Weinberg1995}.  Gauge invariance tells us that
$|\phi|^2$ and $|\phi|^4$ are admissible terms; it is silent on $\mu^2$ and
$\lambda$.  This is not a new principle, and we do not present it as one.

What is worth stating precisely is its consequence \emph{within} the
division-algebra programme.  The programme's appeal is that representation
content and charge structure emerge from the algebra with little external input,
and it is natural to ask how far that reach extends.  Theorem~\ref{thm:deriv}
marks the boundary exactly: the reach extends to the symmetry, the
representations, and the charges, and stops at the dynamics that fix the scale of
breaking.  The value of the theorem is as a boundary marker, made precise in
this setting and engaged against a specific prior claim in the next subsection.
It is not a central discovery of the paper, and a reader who takes it for one has
misread our intent.

\subsection{Relation to Stoica: Clarification, Not Refutation}
\label{sec:deriv-stoica}

Stoica~\cite{Stoica2018}, in the same journal this series targets, constructs
the electroweak sector in the $\cl(6)$ language and states that the electroweak
symmetry is ``already broken by the geometry of the algebra,'' reading off a bare
Weinberg angle $\sin^2\theta_W=0.25$.  Theorem~\ref{thm:deriv} says the algebra
does not fix the Higgs potential.  These two statements are sometimes read as
being in tension.  They are not, and the distinction is instructive.

Stoica's $\sin^2\theta_W=0.25$ is a tree-level generator-normalization ratio.
It reflects how the $\UU(1)_Y$ generator is normalized relative to the
$\su(2)_L$ generators within the algebraic embedding---a ratio of Cartan
normalizations, fixed by the representation content.  It is \emph{not} a
prediction of the potential parameters $\mu^2$ and $\lambda$, and it is not
derived from any dynamics of symmetry breaking.  Theorem~\ref{thm:deriv}
therefore does not refute it: the two statements concern different objects.  The
Weinberg-angle ratio is exactly the kind of quantity the algebra \emph{does} fix
(a normalization built into the embedding); the potential is exactly the kind it
does \emph{not}.

What Theorem~\ref{thm:deriv} does is clarify the phrase ``broken by geometry.''
The geometry of the algebra fixes the symmetry group, the representation
content, and the embedding ratios that yield a tree-level Weinberg angle.  It
does not fix the electroweak \emph{breaking} itself: the sign of $\mu^2$ that
selects the broken vacuum, and the scale $v$ at which breaking occurs, are not
determined by the geometry.  ``Broken by geometry'' is apt for the pattern of
symmetry breaking---which fields can acquire vacuum expectation values, and into
which unbroken subgroup the group descends---but it overstates the reach if read
as fixing the breaking dynamics.  The precise division is: geometry fixes the
pattern and the representation-theoretic ratios; the potential that triggers and
scales the breaking is separate input.  We offer this as a refinement of
Stoica's framing, not a contradiction of his result.

\subsection{What the Algebra Does Fix: Operator Forms}
\label{sec:deriv-yukawa}

The Derivability Limit is precise, not total.  It is worth recording the
complementary positive statement, because it shows the boundary is sharp rather
than a blanket disclaimer.

After symmetry breaking, fermion masses arise from Yukawa couplings of the form
\begin{equation}
	\mathcal{L}_{\mathrm{Yukawa}}
	= -\,y_f\,\bar\psi_L\,\phi\,\psi_R + \text{h.c.}
	\label{eq:yuk}
\end{equation}
The \emph{form} of this coupling is algebraically forced:
$\bar\psi_L\phi\psi_R$ is the unique renormalizable gauge-invariant bilinear
contracting a left-handed doublet, the Higgs doublet, and a right-handed
singlet, with all charges summing to zero.  Gauge invariance, which the algebra
supplies, permits no other renormalizable Yukawa structure.  What the algebra
does not fix is the \emph{magnitude} $y_f$.  The Standard Model carries thirteen
independent Yukawa couplings, spanning six orders of magnitude from the electron
to the top quark, and all thirteen are free parameters from the algebraic point
of view.

This is the same form-versus-coefficient distinction as
Theorem~\ref{thm:deriv}, now on the Yukawa sector: the algebra determines which
couplings may appear and forbids the rest, but assigns none of their values.
The limit is a clean line, not an absence of content.

\subsection{Remark: The One Speculative Handle}
\label{sec:deriv-bankscasher}

\begin{remark}
	\label{rem:bankscasher}
	There is exactly one route by which a dynamical mechanism, rather than pure
	representation theory, might constrain $\mu^2$ within the programme.  If the
	Higgs were a composite object---a fermion condensate, as in technicolor,
	rather than an elementary scalar---then a Banks--Casher-type
	relation~\cite{BanksCasher1980} connecting a non-zero condensate to the
	spectral density of a fundamental operator at zero eigenvalue could in
	principle constrain the breaking, and the programme's spectral-projector
	methods would be well suited to posing the question.  This is the right
	question if the Higgs is composite and the wrong one if it is elementary;
	the discovery of a Higgs boson with properties consistent with an
	elementary scalar disfavors it.  We flag this as the sole caveat that could
	alter the verdict of Theorem~\ref{thm:deriv} and do not pursue it.
\end{remark}

% ============================================================
\section{Discussion and Outlook}
\label{sec:discussion}
% ============================================================

\subsection{What the Algebraic Machinery Determines}
\label{sec:can-do}

Taken together, the results of this paper draw a sharp line around the algebraic
programme.  On one side lie the quantities the algebra fixes.  It fixes the
representation content: which irreducible representations of the Standard Model
gauge group appear, recovered for the color sector from the $\cl(6)$ ideal of
Papers~1 and~2 and for one full generation from the $\cl(10)$ construction of
Sec.~\ref{sec:cl10}.  It fixes the charge assignments, with the
Gell-Mann--Nishijima relation $Q=T_3+Y$ satisfied by all sixteen states of a
generation.  It fixes the pattern of symmetry breaking once a vacuum direction
is chosen, and the embedding ratios that yield a tree-level Weinberg angle, in
the precise sense clarified in Sec.~\ref{sec:deriv-stoica}.  It fixes the form of
the allowed interactions: the unique gauge-invariant Yukawa bilinear and quartic
scalar terms, as set out in Sec.~\ref{sec:deriv-yukawa}.  And it places the
construction at a definite algebraic level: $\cl(6)$ is maximal for QCD
(Corollary~\ref{cor:maximal}), and $\cl(10)=\cl(6)\grtensor\cl(4)$ is the minimal
graded-tensor extension that restores the weak-doublet assignment
Theorem~\ref{thm:schur} forbids within $\cl(6)$ alone.

A separate, more modest positive should be recorded for what it is.  The
per-configuration completeness identity $\sum_r C^{(r)}=C$, satisfied to order
$10^{-16}$ on every gauge configuration, is a genuine consistency check on the
propagator and projector code.  It is a diagnostic, not a physical result, and
we claim nothing more for it.

\subsection{What It Does Not Determine}
\label{sec:cannot-do}

On the other side lie the no-go results.  The algebra does not fix the Higgs
potential---neither the sign of $\mu^2$ nor the magnitude of $\lambda$
(Theorem~\ref{thm:deriv})---and therefore fixes neither the electroweak scale
$v$ (Corollary~\ref{cor:vscale}) nor the Higgs mass $m_H$
(Corollary~\ref{cor:mhmass}).  It does not fix the thirteen Yukawa magnitudes or
the fermion mass hierarchy they encode (Sec.~\ref{sec:deriv-yukawa}).  The
diagonal electroweak grade channel carries no dynamics at all: its grade
fractions are pinned to one-half by construction, at every coupling
(Proposition~\ref{prop:rigidity}).  The algebra does not select left over right:
both $\su(2)_L$ and $\su(2)_R$ are present in the Pati--Salam structure of
Sec.~\ref{sec:patisalam}, and gauging only the left copy---the parity violation
of the physical world---is an external input the algebra accommodates but does
not supply.  And, as Theorem~\ref{thm:schur} establishes, no color-commuting
weak isospin lives inside $\cl(6)$ in the first place.

\subsection{The Bridge Between Three Programmes}
\label{sec:bridge}

Beyond delimiting the machinery, the construction of Sec.~\ref{sec:cl10} serves a
connective purpose.  The explicit graded-tensor realization
$\cl(10)=\cl(6)\grtensor\cl(4)$ takes the lattice-ready $\cl(6)$ framework of
this series as a literal tensor factor, and in doing so links three otherwise
separate lines of work: the lattice $\cl(6)$ construction developed here, the
algebraic Standard-Model programme of Furey~\cite{Furey2018a,Furey2018b} and
Gresnigt~\cite{Gresnigt2020}, and Catterall's K\"ahler--Dirac route to
Pati--Salam on the lattice~\cite{Catterall2023,Catterall2024}.  This bridge is the paper's
positive role alongside its two no-go results, and it locates the lattice
programme within the broader algebraic landscape rather than apart from it.

\subsection{Outlook}
\label{sec:outlook}

The line drawn here is not the end of the programme but its honest boundary.  The
diagonal grade channel is inert, but the off-diagonal cross-channel
$\Delta_{\mathrm{cross}}$ introduced in Sec.~\ref{sec:rigidity} is
coupling-sensitive and is the natural carrier of gauge dynamics in this
decomposition; its analysis, including the gauge-fixing systematics it entails,
is taken up elsewhere.  That channel, and the composite-Higgs question flagged in
Remark~\ref{rem:bankscasher}, are the two directions in which the boundaries
established in this paper might eventually be pushed.

% ============================================================
\section{Conclusions}
\label{sec:conclusions}
% ============================================================

This paper has drawn a precise boundary around what the $\cl(6)$-based algebraic
programme determines, through three results.

First, a Schur obstruction (Theorem~\ref{thm:schur}): no weak isospin commuting
with color can act on the eight-dimensional $\cl(6)$ ideal except on the color
singlets, so quarks cannot be weak doublets within $\cl(6)$.  The obstruction is
exact and embedding-independent, and it establishes that $\cl(6)$ is
algebraically maximal for QCD (Corollary~\ref{cor:maximal}).

Second, the minimal resolution (Sec.~\ref{sec:cl10}): the graded tensor product
$\cl(10)=\cl(6)\grtensor\cl(4)$ adjoins a $\cl(4)$ factor on which weak isospin
acts, restoring the doublet assignment via the Pati--Salam embedding and
reproducing all sixteen charges of one generation.  This construction is a
structural re-expression of known Pati--Salam content, and its value is as an
explicit bridge among the lattice $\cl(6)$, Furey--Gresnigt, and Catterall
programmes.

Third, two no-go statements bounding the dynamics: the diagonal electroweak
grade decomposition is inert, equal to one-half at all couplings
(Proposition~\ref{prop:rigidity}); and the algebra fixes no parameter of the
Higgs potential, hence neither the electroweak scale nor the Higgs mass
(Theorem~\ref{thm:deriv}).  Together they mark, precisely, where algebraic
structure ends and dynamics begins.

% ============================================================
\begin{thebibliography}{99}
% ============================================================

	%% Ordered by first citation.  Plain \bibitem{key} numbers sequentially.

	%% Papers in this series
	% NOTE: Replace journal name, arXiv ID, and year for P1/P2 before submission.
	\bibitem{FirestoneP1}
	T.~M. Firestone,
	``Clifford algebra $\cl(6)$ formulation of lattice QCD: algebraic
	structure, Dirac operator construction, and verification against standard
	benchmarks,''
	[arXiv:XXXX.XXXXX] ([year]).

	\bibitem{FirestoneP2}
	T.~M. Firestone,
	``Grade decomposition of the lattice QCD pion correlator in the $\cl(6)$
	framework,''
	\textit{Adv.\ Appl.\ Clifford Algebras} (submitted, [year]).

	%% Furey programme (cited as a block in §1)
	\bibitem{Furey2012}
	C.~Furey,
	``Unified theory of ideals,''
	\textit{Phys.\ Rev.}\ D~\textbf{86}, 025024 (2012).

	\bibitem{Furey2014}
	C.~Furey,
	``Generations: three prints, in colour,''
	\textit{JHEP}~\textbf{2014}, 046 (2014).

	\bibitem{Furey2015}
	C.~Furey,
	``Standard model physics from an algebra?''
	Ph.D. thesis, University of Waterloo (2015); arXiv:1611.09182.

	\bibitem{Furey2018a}
	C.~Furey,
	``Three generations, two unbroken gauge symmetries, and one
	eight-dimensional algebra,''
	\textit{Phys.\ Lett.}\ B~\textbf{785}, 84 (2018).

	\bibitem{Furey2018b}
	C.~Furey,
	``$\su(3)_C\times\su(2)_L\times\UU(1)_Y\,(\times\,\UU(1)_X)$ as a symmetry
	of division algebraic ladder operators,''
	\textit{Eur.\ Phys.\ J.}\ C~\textbf{78}, 375 (2018).

	\bibitem{Stoica2018}
	O.~C. Stoica,
	``Leptons, quarks, and gauge from $C\ell_6$,''
	\textit{Adv.\ Appl.\ Clifford Algebras}~\textbf{28}, 52 (2018).

	%% Schur's lemma (§2)
	\bibitem{FultonHarris}
	W.~Fulton and J.~Harris,
	\textit{Representation Theory: A First Course},
	Graduate Texts in Mathematics~\textbf{129} (Springer, New York, 1991).

	%% Cl(6)(x)Cl(4)=Cl(10) factorization (§3)
	\bibitem{Gresnigt2020}
	N.~G. Gresnigt,
	``The Standard Model particle content with complete gauge symmetries from
	the minimal ideals of two Clifford algebras,''
	\textit{Eur.\ Phys.\ J.}\ C~\textbf{80}, 583 (2020);
	arXiv:2003.08814.
	\href{https://doi.org/10.1140/epjc/s10052-020-8141-1}{doi:10.1140/epjc/s10052-020-8141-1}.

	%% Pati-Salam embedding and branching (§3) --- co-cited, kept adjacent
	\bibitem{PatiSalam1974}
	J.~C. Pati and A.~Salam,
	``Lepton number as the fourth color,''
	\textit{Phys.\ Rev.}\ D~\textbf{10}, 275 (1974).

	\bibitem{BaezHuerta2010}
	J.~C. Baez and J.~Huerta,
	``The algebra of grand unified theories,''
	\textit{Bull.\ Amer.\ Math.\ Soc.}~\textbf{47}, 483 (2010);
	arXiv:0904.1556.
	\href{https://doi.org/10.1090/S0273-0979-10-01294-2}{doi:10.1090/S0273-0979-10-01294-2}.

	%% Catterall bridge (§3) --- co-cited, kept adjacent.
	% PRD 107, 014501 (2023) Sec. VII: Pati-Salam from reduced Kahler-Dirac.
	% SciPost Phys. 16, 108 (2024): chiral-fermion connection.
	\bibitem{Catterall2023}
	S.~Catterall,
	``'t~Hooft anomalies for staggered fermions,''
	\textit{Phys.\ Rev.}\ D~\textbf{107}, 014501 (2023).
	\href{https://doi.org/10.1103/PhysRevD.107.014501}{doi:10.1103/PhysRevD.107.014501}.

	\bibitem{Catterall2024}
	S.~Catterall,
	``Lattice regularization of reduced K\"ahler--Dirac fermions and connections
	to chiral fermions,''
	\textit{SciPost Phys.}~\textbf{16}, 108 (2024).
	\href{https://doi.org/10.21468/SciPostPhys.16.4.108}{doi:10.21468/SciPostPhys.16.4.108}.

	%% EFT form-vs-coefficient lore (§5.2)
	\bibitem{Weinberg1995}
	S.~Weinberg,
	\textit{The Quantum Theory of Fields, Vol.~1: Foundations}
	(Cambridge University Press, 1995).

	%% Banks-Casher remark (§5.5)
	\bibitem{BanksCasher1980}
	T.~Banks and A.~Casher,
	``Chiral symmetry breaking in confining theories,''
	\textit{Nucl.\ Phys.}\ B~\textbf{169}, 103 (1980).

\end{thebibliography}

\end{document}
